\section{Good-orbit counts and uniform Type-II data}
\label{sec:type-two-family}

\status{We prove the actual family Type-II hypotheses and their local
application. Separately, a finite probability argument produces a subset
with small weighted bad-orbit mass. Identifying that subset with the one
used by the arithmetic host remains an explicit construction step.}

Fix $0<\epsilon<1/52$, $\zeta=\epsilon/J>0$, and the nonnegative global
constant tuple $\mathbf C$. The literal parameters are
\[
\begin{gathered}
 \delta=x^{-\epsilon/8},\quad h=2+8\epsilon,\quad
 K_0=\lceil8(h+2)/\epsilon\rceil,\quad r=K_0+1,\\
 u=\lceil\delta^{-1}\rceil,\quad \ell=(6r+1)u,\quad
 P_0=10^{-(6r+1)},\quad M=\lceil(h+3)\log x/P_0\rceil,\\
 G=(M+5)\ell,\qquad
 N_1=\lceil\max(6eG\delta,(h+3)\log_2x)\rceil+1.
\end{gathered}\tag{GII.1}
\]
Thus $G=\texttt{Gp}$ and $N_1=\texttt{N1}$, including their ceilings.
All assertions below use one sufficiently large $x$ threshold after the
fixed constants have been chosen.

\subsection{The finite counting input}

\paragraph{GII1: enough tops inside each good orbit.}
Let $T\subset\mathbb Z/n'\mathbb Z$ satisfy the actual good-orbit
conditions, with $3G\le n'$. A marked complete tile of length $\ell$
has one selected even point in each even-numbered cell of length $u$,
and no selected points in odd-numbered cells. Every $M+1$ successive
complete tiles contain a marked one. List the marked tile starts
periodically on $\mathbb Z$. Their successive gaps are less than
$(M+2)\ell$: the wrap contributes the remainder of $n'$ upon division
by $\ell$, which is less than $\ell$. From any position, the next
marked tile's second cell therefore supplies an even top at positive
distance at most $G$.

The map $(t,d)\mapsto t-d$, for $t\in T$ and $1\le d\le G$,
covers the orbit. Conversely, double-counting all cyclic $G$-windows,
each of which contains at most $N_1$ tops, gives
\[
 n'\le |T|G\le n'N_1.\tag{GII.2}
\]
The upper bound improves the source's $|T|G\le N_1(n'+G)$; we retain
its downstream constants. The least positive even successor has gap
$\le G$ and at most $N_1$ interior tops. It is distinct because $G<n'$.
Hence there are at least two even tops, and their literal cyclic
segments satisfy $\texttt{segG}\le G$ and $\texttt{segS}\le N_1$.
This proves \texttt{OrbitSegs}. When $L_1\le L_2$ and the short-window
cap is $r-1$, fewer than $r$ even tops occur in the preceding $L_1$
positions. These are the remaining positional clauses of
\texttt{AuxBlock8.adm1}.

\paragraph{GII2--GII3: one uniform constant and the actual orbit map.}
For $x\ge3$, put
\[
 C_G=2(6r+1)\bigl((h+3)10^{6r+1}+6\bigr),\qquad C_e=20C_G.
\]
The inequalities $1\le u\delta\le2$, $M< (h+3)\log x/P_0+1$ and
$\log x\ge1$ give $G\delta\le C_G\log x$. Moreover $G\delta\ge5$ and
$(h+3)\log_2x\le2M\le2G\delta$, so the first argument of the maximum
in (GII.1) dominates. Thus $N_1<6eG\delta+2\le16.8G\delta$. Consequently
\[
 G\le\frac{C_e\log x}{\delta},\quad N_1\le C_e\log x,\qquad
 \frac{n'\delta}{C_e\log x}\le N:=|T|\le40n'\delta.
 \tag{GII.3}
\]
For an inert prime, $y\mapsto t_0+yh_p$ is the actual orbit bijection,
and the representative identities identify $N$ with \texttt{Nv}.
The predicate \texttt{IsV} means that the unique index $y$ is even.
Transporting the preceding counts and distances proves
\texttt{adm1} for every $C_{31}\ge C_e$.

In a supplied family, $r,h,C_e$ depend only on $\epsilon$. For
$x\ge\max(3,x_1+1)$ these statements hold for every member. If
$\mathbf C.\mathrm{CN}\ge C_e$, they also prove the exact
\texttt{FamilyWindowLogBound}. Nonnegative substitution gives
$\rho\le\rho_B\le x^{-2}$, using the family's stated bound on $\rho_B$.
No current family index is selected; empty index sets cause no difficulty.

\subsection{Supplying every Type-II inequality}

Write
\[
\begin{gathered}
 a=xe^{x^\epsilon},\quad X'=e^{x^{\epsilon/3}},\quad b'=a/X',
 \quad p_0=k^2x^{26},\quad Q=p^G,\\
 L_1(p)=\lceil\log_p((kx^{26})^4)\rceil,\qquad
 L_2=2\sup_{p\ \mathrm{inert}}L_1(p),\quad c=4\mathrm{Sc},\quad c_0=2b'+2c.
\end{gathered}\tag{GII.4}
\]
The relevant \texttt{SizesAt} clauses give
$x^{9/2}\le k\le x^5\le x^6$, $n\le x^3$ and $\mathrm{Sc}\le b'/x$.
The actual parameter-fill record gives
$\Lambda:=C_{\rm La}N_1^2\log a=\mathrm{LaG}$.
Eventually $c\le b'$, $c_0\le5b'$ and $b'\ge1$.

\paragraph{GII4: levels, prime size and weights.}
Since $\log(kx^{26})\le31\log x$ and $\log p\ge\log2>0.69$,
$L_2\le514\log x<u$ eventually; its definition gives $2L_1\le L_2$.
Also $3G\le x^{\epsilon/4}$ eventually. The host's small primes split
up to $\lfloor x^{\zeta/5}\rfloor$; inert classification excludes them
and $2$. Hence $p>x^{\zeta/5}$, so $p\ge13$ eventually, and
\[
 L_2\le L_b:=2\lceil620/\zeta\rceil.\tag{GII.5}
\]
The latter constant will also control the probability argument.

Let $I_b=2n\delta x^\epsilon\Lambda b'/k$. If $\rho\le1/7$, its
nonnegative summands give $I_b\le b'$. The actual integer weights lie in
$[a,a+a/(kx^{13})]$ and $[b',b'+I_b]$, precisely the global windows.
For $x\ge24$, $c\le b'/6$ and $2b'\le a$. Therefore, with
$c''=\lfloor c\rfloor$,
\[
 (a-c)+c''\le a_t\le2(a-c)-c'',\qquad
 b_t+2c''\le\lfloor c_0\rfloor.\tag{GII.6}
\]
For the upper $a_t$ bound, use $a_t\le3a/2$ and $3c\le a/4$.
Thus all shifted and rounded weight clauses are supplied.

\paragraph{GII5: the lower count supplies the rank margin.}
Assume $p<p_0$, $x^{\epsilon/4}\le n'\le n$ and $N\le n'$.
Put $y=x^{\epsilon/24}$. Choose $x$ so that both $C_e\log x$ and
$C_{31}\log x$ are at most $y$. The lower count gives
$Ny\ge n'\delta\ge y^3$, whence
\[
 N_1^2\le y^2\le N\le rN.\tag{GII.7}
\]
Also $r,N_1\le x$ eventually, so $N_1+3r\le4x\le k$ and
$r\le x^{26}$. Since $\log a\le2x$,
$C_k'N_1^2\log a\le2C_k'x^3\le k$ eventually. Finally the ceiling in
$L_1$ gives $p^{L_1}\ge(kx^{26})^4$. Dividing by $p<p_0$ yields
\[
 256k^2N(x^{26})^2\le k^2x^{78}\le p^{L_1-1},\tag{GII.8}
\]
using $N\le x^3$ and $x\ge500$.

\paragraph{GII6: the full exponential size budget.}
Let $N_+=40n'\delta$ and use the actual denominator
\[
 D=\left\lceil C_DN_+(N_1^3+r^5+k^2)\log_pa
                 +2rn'c+rN_+c_0\right\rceil.\tag{GII.9}
\]
Because $p\le x^{38}$, (GII.3) gives
$\log Q\le38C_G(\log x)^2x^{\epsilon/8}\le x^{7\epsilon/48}$
eventually. Put $E=\exp(5x^{7\epsilon/48})$, so $Q^5\le E$.
For $x\ge\max(500,480C_D,C_s)$, $\log_pa\le4x$ and
\[
\begin{aligned}
 D&\le x^{17}+202x^4b'+1\le x^{18}b',\\
 8(D+1)+rNc_0+2rn'c&\le9x^{18}b',\\
 C_sk^3r^2NQ^5[8(D+1)+rNc_0+2rn'c]x^{26}
   &\le9x^{68}Eb'\le x^{69}Eb'.
\end{aligned}\tag{GII.10}
\]
These estimates include the ceiling's extra unit. Since
$7\epsilon/48<\epsilon/3$, choose $x$ so that
$70\log x+5x^{7\epsilon/48}\le x^{\epsilon/3}$. Then
$x^{69}E+1\le x^{70}E\le X'$, and the last line of (GII.10) is
$\le a-b'\le a-c$. This proves every field of \texttt{NumII}.

\paragraph{GII7--GII8: the literal family application.}
A candidate prime has $n'=\operatorname{ord}(h_p)\ge x^{\epsilon/4}$;
its order divides $n$. For each good representative, (GII.3) supplies
the count used above. We have now proved all of \texttt{HypIIAt},
uniformly in the prime, orbit and family member. In a \texttt{FamP},
\texttt{sizes}, \texttt{fills} and the bound on $\rho_B$ supply exactly
the records used here. Thus
\[
 \exists x_0\ \forall x\ge x_0\ \forall i\in\mathrm{Ix}(x),
 \qquad\texttt{HypIIAt}(D_{x,i},L_{x,i},X_{x,i},\mathbf C).
 \tag{GII.11}
\]
The constant tuple and $\epsilon,J$ precede $x_0$; $i,p$ and the orbit
follow it. If no family with the given $\epsilon$ exists the universal
claim is vacuous, rather than requiring an invalid division by $\epsilon$.

For completeness, use $c\ge0$ and $a-c\ge1$, already proved above.
Take a local constant $C\ge4129$ and require
$C\le C_k,C_t,C_{\rm La},C_D,C_s,C_k'$. Apply EA9 at
\[
 m=k-1,\quad a_0=a-c,\quad Y=x^{26},\quad
 c''=\lfloor c\rfloor,\quad c_{0,\mathrm{int}}=\lfloor c_0\rfloor,
 \quad D_p=D,\quad c_*=8(D+1).
\]
Equations (GII.6--10) verify its hypotheses. In particular
$CNk^2\log_pa_0\le D$. Monotonicity of every nonnegative height term
gives the exact global height
\[
 t_{\rm II}=C_t rN\left(\frac{\Lambda a}{pk}+ap^{-L_1}+c_0+kQ\right)+2rn'c.
\]
Hence the local witness works at $\lceil t_{\rm II}\rceil$ and loss
$\kappa_{\rm II}=\min(1,C_kNr^2kx^{-26})$. The truncation at $1$ uses
the empty witness when its rank requirement is nonpositive. This proves
the actual \texttt{adm6}, at the supplied standard coefficient setting;
the finite extension to an actual place is separate.

\subsection{Probability and a simultaneous choice of tops}

\paragraph{GII9: condition on a selected point correctly.}
Select the ambient points independently with probability $\delta$ and
condition on a fixed point being selected. Its orbit has one forced
index $i_0$; the other indices remain independent. At most $L_2$ short
windows contain $i_0$. Failure of their cap $K_0$ therefore has probability
at most
\[
 L_2\binom{L_2}{K_0}\delta^{K_0}
 +n'\binom{L_2}{K_0+1}\delta^{K_0+1}.
\]
Since $\delta^{K_0}\le x^{-(h+2)}$, $n'\le n\le x^{h+\epsilon/8}$
and $L_2\le L_b$, this is at most $(L_b+1)2^{L_b}x^{-2}$.
The additional $\delta$ in the second term is essential.

For $\delta\le1/10$, $(1-\delta)^u\ge1/4$ and
$\lfloor u/2\rfloor\delta\ge9/20$. Thus an empty cell has probability
$\ge1/4$, and exactly one selected even point has probability $\ge9/80$.
A complete tile is marked with probability at least $P_0$.
Among any $M+1$ distinct consecutive tiles, discard the one containing
$i_0$, if any. The other $M$ marking events are independent, so a union
bound gives tile failure at most
$n'(1-P_0)^M\le n'x^{-(h+3)}\le x^{-2}$.

Put $B=N_1-1$. For a long window, the conditional failure probability
is bounded by the binomial tail
\[
 \binom GB\delta^B\le(eG\delta/B)^B\le2^{-B}\le x^{-(h+3)}.
\]
If $B>G$ the tail is zero. A union bound over $n'$ windows contributes
at most $x^{-2}$. Therefore
\[
 \Pr(\text{bad orbit}\mid\text{point selected})\le C_bx^{-2},
 \qquad C_b=(L_b+1)2^{L_b}+2.\tag{GII.12}
\]
Marginalizing over points outside the orbit multiplies by $1$; removing
the conditioning multiplies by $\delta$. This is the literal
\texttt{BadProb}, uniform over all sufficiently long steps and all points.

\paragraph{GII10: one subset works for every prime in the finite sum.}
Let $\mu=\delta|H|\ge16$. Chebyshev gives size between $\mu/2$ and
$2\mu$ with probability at least $3/4$. For nonnegative weights
$\omega_p$, the weighted bad count $Z$ has expectation at most
$B_0=C_bx^{-2}\mu\sum_p\omega_p$. Markov gives $Z\le4B_0$ with
probability at least $3/4$; if $B_0=0$, then $Z=0$ almost surely.
The events intersect, proving the exact existential \texttt{TopsOK}.

Take $\omega_p=\log p/p$ over the source's primes $p<X'$ of large
orbit order. For $n=\lfloor y\rfloor\ge1$,
\[
 \frac n2\sum_{p\le n}\frac{\log p}{p}
 \le\sum_{p\le n}\lfloor n/p\rfloor\log p
 \le\log(n!)\le n\log n.
\]
Hence $\sum_{p<y}\log p/p\le2\log y$, including the empty small cases.
The selected subset consequently satisfies
\[
 Z\le16C_b|T|x^{-2+\epsilon/3}
   \le |T|x^{-2+\epsilon}\tag{GII.13}
\]
eventually. This is a simultaneous weighted abundance statement. If this
subset is the actual host subset, it is exactly \texttt{adm2}; for an
arbitrary supplied host, that identification remains to be proved.
